\documentclass[12pt]{article}
\usepackage[utf8]{inputenc}
\usepackage[T1]{fontenc}
\usepackage{hyperref}
\title{Iten Language Resources for K-12 Curriculum Development}
\author{Compiled by Claude Sonnet 5.5 (claude-sonnet-5-5) for Andr\'as Kornai}
\date{}
\begin{document}
\maketitle
\section{Introduction}
Iten (Eten, Tyen; ISO 639-3 etx, Glottocode eten1239) is a Plateau (Benue-Congo) language of central Nigeria. Sources place it in Plateau and Kaduna States and disagree on speaker numbers (40,000 to about 75,000). Little educational material was found.
\section{Primary Educational Resources}
None were found in the search results.
\section{Linguistic and Reference Materials}
\begin{itemize}
\item Wikipedia: \url{https://en.wikipedia.org/wiki/Eten_language}
\item Glottolog (classification): \url{https://glottolog.org/resource/languoid/id/eten1239.bigmap.html}
\item Numeral system notes from a 1994 survey: \url{https://lingweb.eva.mpg.de/channumerals/Iten.htm}
\item ASJP word list entry: \url{https://asjp.clld.org/languages/ITEN}
\end{itemize}
\section{Religious and Community Resources}
Global Recordings Network lists Bible-related audio: \url{https://globalrecordings.net/en/language/etx}
\section{Digital and Audio Resources}
Wikitongues language page: \url{https://wikitongues.org/languages/etx/}
\section{Curriculum Development Context}
No school use is known. Speakers often also understand Hausa and Birom, which matters for any materials plan.
\end{document}